% !TeX document-id = {3d5ac9d1-fe14-4307-8745-89369ff1b702}
% !TeX TXS-program:compile = txs:///latexmk/{-pdf}
% !TeX root = main.tex
% Kompilacja: latexmk main.tex
% Windows: latexmk wymaga MiKTeX-a oraz interpretera Perl, np. Strawberry Perl.
% Prism/Overleaf: wybierz main.tex i kompilator LuaLaTeX; lokalny Perl nie jest potrzebny.

% ============================================================================
% JAK KORZYSTAĆ Z TEGO PLIKU
% ============================================================================
% Ten plik określa kolejność elementów całego dokumentu. Nie wpisuj tutaj
% treści raportu ani danych autorów.
%
% MOŻNA ZMIENIAĆ:
%   - kolejność poleceń \include w sekcji „ROZDZIAŁY PROJEKTU”;
%   - listę dołączanych rozdziałów i załączników;
%   - tymczasowo wyłączać rozdział, dodając przed \include znak %.
%
% NIE NALEŻY ZMIENIAĆ:
%   - poleceń między \begin{document} a sekcją „ROZDZIAŁY PROJEKTU”;
%   - poleceń \input konfiguracji, \documentclass i \usepackage poniżej;
%   - poleceń \begin{document}, \end{document} i \appendix.
%
% Dane projektu uzupełniaj w: config/metadata.tex
% Treść raportu zapisuj w: chapters/*.tex
% Wygląd dokumentu definiuje: config/project-style.sty
% ============================================================================

% === KONFIGURACJA DOKUMENTU — NIE ZMIENIAJ TYCH TRZECH POLECEŃ ===
% ============================================================================
% DANE PROJEKTU — TEN PLIK NALEŻY UZUPEŁNIĆ
% ============================================================================
% Zmieniaj wyłącznie tekst znajdujący się między drugą parą nawiasów { }.
% Nie usuwaj poleceń \newcommand ani ich nawiasów.
%
% Przykład:
%   \newcommand{\projecttitle}{Projekt stanowiska pomiarowego}
%
% Znaki specjalne LaTeX-a, m.in. %, &, _ oraz #, należy poprzedzić znakiem \.
% Przykład: 50\%, R\&D, plik\_wynikowy
% ============================================================================

% ETAP PRACY: draft albo final.
% draft — brakujące dane powodują ostrzeżenia, ale PDF nadal powstaje.
% final — brakujące podstawowe dane zatrzymują kompilację.
\newcommand{\projectstatus}{draft}

% === UCZELNIA I PODSTAWOWE DANE PROJEKTU ===
\newcommand{\university}{Politechnika Opolska}
\newcommand{\faculty}{Wydział Elektrotechniki, Automatyki i Informatyki}
\newcommand{\projecttitle}{Tytuł projektu}
% Podtytuł jest opcjonalny. Pozostaw {}, jeżeli nie jest potrzebny.
\newcommand{\projectsubtitle}{}
\newcommand{\coursetitle}{Nazwa przedmiotu}
\newcommand{\projecttype}{Projekt semestralny}
\newcommand{\programme}{Nazwa kierunku}
\newcommand{\studyform}{studia stacjonarne / niestacjonarne}
\newcommand{\semester}{semestr 0}
\newcommand{\instructor}{dr inż. Imię i nazwisko}

% === AUTOR LUB ZESPÓŁ ===
% Znak & oddziela nazwisko od numeru albumu, a \\ kończy wiersz.
% Jeden autor:
%   Imię i nazwisko & 000000 \\
% Zespół:
%   Anna Nowak & 123456 \\
%   Jan Kowalski & 234567 \\
\newcommand{\projectauthors}{%
  Imię i nazwisko & 000000 \\
}
% Wpisz tych samych autorów w jednym wierszu. To pole jest używane wyłącznie
% we właściwościach PDF, dlatego nie stosuje się tu & ani \\.
\newcommand{\projectpdfauthors}{Imię i nazwisko}

% === MIEJSCE I TERMIN ===
\newcommand{\city}{Opole}
\newcommand{\academicyear}{2026/2027}
\newcommand{\submissiondate}{dd.mm.rrrr}

% === LOGO ===
% Ścieżka jest liczona od katalogu głównego projektu.
% Logo można wyłączyć przez ustawienie: \newcommand{\universitylogo}{}
\newcommand{\universitylogo}{assets/branding/logo-weaii.jpeg}
\newcommand{\universitylogowidth}{0.55\textwidth}

% === OPCJONALNE WYKAZY ===
% Dozwolone wartości to wyłącznie true albo false, bez cudzysłowów.
% Wykaz ma sens dopiero wtedy, gdy raport zawiera odpowiednie podpisane elementy.
\newcommand{\showlistoffigures}{false}
\newcommand{\showlistoftables}{false}
\newcommand{\showlistoflistings}{false}

% === OŚWIADCZENIA KOŃCOWE ===
% Obie części są domyślnie włączone. Wartość false można zastosować wyłącznie,
% gdy prowadzący nie wymaga danego oświadczenia albo nie ma ono zastosowania.
\newcommand{\showaideclaration}{true}
\newcommand{\showcontributions}{true}

\documentclass[11pt,a4paper,oneside]{report}
\usepackage{config/project-style}

% === POCZĄTEK DOKUMENTU — NIE USUWAJ ===
\begin{document}

% === STRONA TYTUŁOWA I SPISY — USTAWIENIA SĄ W config/metadata.tex ===
% Widoczność spisów rysunków, tabel i listingów ustawia się tam przez true/false.
% Poniższych poleceń nie trzeba zmieniać ani usuwać.
\hypersetup{pageanchor=false}
\pagenumbering{gobble}
\makeprojecttitle

\hypersetup{pageanchor=true}
\pagenumbering{roman}
\begin{singlespace}
  \hypersetup{linkcolor=black}
  \tableofcontents
  \hypersetup{linkcolor=projectblue}
\end{singlespace}
\ifdefstring{\showlistoffigures}{true}{\printprojectlistoffigures}{}
\ifdefstring{\showlistoftables}{true}{\printprojectlistoftables}{}
\ifdefstring{\showlistoflistings}{true}{\printprojectlistoflistings}{}
\clearpage
\pagenumbering{arabic}

% === ROZDZIAŁY PROJEKTU — TĘ LISTĘ MOŻNA ZMIENIAĆ ===
% Każde polecenie \include dołącza jeden plik z katalogu chapters/.
% Nie podawaj rozszerzenia .tex.
%
% Przykład wyłączenia rozdziału:
% % % TREŚĆ ROZDZIAŁU: zastąp wskazówki opisem przebiegu własnego projektu.
% Cały rozdział można wyłączyć w main.tex, jeśli prowadzący go nie wymaga.
\chapter{Organizacja i przebieg projektu}
\label{chap:organizacja}

\section{Etapy i harmonogram}
Porównaj planowane etapy z rzeczywistym przebiegiem. Wyjaśnij istotne zmiany zakresu lub terminów.

\section{Podział zadań}
W projekcie zespołowym jednoznacznie opisz odpowiedzialność poszczególnych osób. W projekcie indywidualnym wskaż własny wkład i wykorzystane elementy zewnętrzne.

\section{Zarządzanie wersjami i materiałami}
Krótko opisz repozytorium, zasady nazywania wersji, strukturę plików oraz sposób przechowywania danych i dokumentacji.

%
% Przykład dodania własnego rozdziału:
% 1. Utwórz plik chapters/08-instrukcja-uzytkownika.tex.
% 2. Rozpocznij go poleceniem \chapter{Instrukcja użytkownika}.
% 3. Dodaj poniżej: \include{chapters/08-instrukcja-uzytkownika}
%
% Kolejność rozdziałów można dostosować do charakteru projektu i wymagań
% prowadzącego.
% TREŚĆ ROZDZIAŁU: usuń teksty instruktażowe i wpisz własne akapity.
% Nie zmieniaj \chapter ani \label, jeśli inne rozdziały odwołują się do tej
% etykiety. Nowe podrozdziały dodawaj przez \section{Tytuł}.
\chapter{Wprowadzenie}
\label{chap:wprowadzenie}

% Usuń wskazówki i zastąp je treścią projektu.
\begin{quote}\itshape
Krótko przedstaw kontekst, potrzebę realizacji projektu i problem, który ma zostać rozwiązany. Wyjaśnij znaczenie tematu dla odbiorcy. Nie opisuj tu jeszcze szczegółowo wymagań ani wyników.
\end{quote}

\section{Kontekst i motywacja}
Opisz dziedzinę, użytkownika lub proces, którego dotyczy projekt. Twierdzenia techniczne i dane pochodzące z literatury poprzyj cytowaniami, np. \cite{lamport1994}.

\section{Problem projektowy}
Zdefiniuj rzeczywistą trudność, a nie tylko czynność do wykonania. Wskaż stan początkowy i pożądany rezultat.

\section{Struktura dokumentu}
Jednym akapitem wyjaśnij rolę kolejnych rozdziałów raportu.

% TREŚĆ ROZDZIAŁU: zastąp wskazówki opisem własnego projektu.
% Zachowaj identyfikatory wymagań, aby można było odwołać się do nich w testach.
% Przykładową tabelę można zmienić albo usunąć w całości.
\chapter{Cel, zakres i wymagania}
\label{chap:cel-wymagania}

\section{Cel projektu}
Sformułuj jeden główny, sprawdzalny rezultat projektu oraz ewentualne cele szczegółowe.

\section{Zakres}
Wskaż funkcje, elementy i etapy objęte projektem. Jawnie zapisz istotne elementy pozostające poza zakresem.

\section{Założenia i ograniczenia}
Opisz ograniczenia techniczne, czasowe, sprzętowe, programowe i organizacyjne.

\section{Wymagania}
Nadaj wymaganiom stabilne identyfikatory i kryteria akceptacji. Przykład przedstawiono w \cref{tab:wymagania}.

\begin{table}[htbp]
  \centering
  \caption{Przykładowy rejestr wymagań}
  \label{tab:wymagania}
  \begin{tabularx}{\textwidth}{@{}p{1.5cm}X X@{}}
    \toprule
    ID & Wymaganie & Kryterium akceptacji \\
    \midrule
    F-01 & System realizuje wskazaną funkcję. & Test T-01 zakończony wynikiem pozytywnym. \\
    NF-01 & Czas odpowiedzi nie przekracza ustalonej wartości. & Pomiar w zdefiniowanych warunkach. \\
    \bottomrule
  \end{tabularx}
\end{table}

% TREŚĆ ROZDZIAŁU: zastąp wskazówki opisem własnej koncepcji i projektu.
% Rysunki przechowuj w katalogu figures/ rozdziału, podpisuj przez \caption
% i oznaczaj przez \label, aby później przywołać je poleceniem \cref.
\chapter{Koncepcja i projekt rozwiązania}
\label{chap:projekt}

\section{Analiza wariantów}
Przedstaw rozważane rozwiązania, kryteria porównania i uzasadnienie wyboru. Nie ograniczaj uzasadnienia do stwierdzenia, że wariant był najłatwiejszy.

\section{Architektura rozwiązania}
Pokaż moduły, granicę systemu, przepływ informacji lub energii oraz interfejsy. Każdy schemat omów w tekście.

\section{Projekt szczegółowy}
Opisz algorytmy, modele, struktury danych, interfejsy, podzespoły, schematy lub parametry niezbędne do wykonania rozwiązania. Liczby z jednostkami zapisuj np. jako \qty{230}{\volt}.

\section{Plan weryfikacji}
Powiąż przypadki testowe z wymaganiami. Dla każdego testu określ dane wejściowe, procedurę, oczekiwany rezultat i warunek zaliczenia.

% TREŚĆ ROZDZIAŁU: zastąp wskazówki opisem wykonania rozwiązania.
% Nie wklejaj długiego kodu bezpośrednio do pliku .tex. Umieść go w katalogu
% code/ rozdziału i wczytaj przez \lstinputlisting — przykład jest na końcu.
\chapter{Realizacja projektu}
\label{chap:realizacja}

\section{Środowisko i narzędzia}
Wymień tylko te wersje narzędzi, bibliotek, urządzeń i konfiguracji, które są istotne dla odtworzenia projektu.

\section{Implementacja lub wykonanie}
Opisz najważniejsze etapy budowy rozwiązania i decyzje podjęte podczas realizacji. Oddziel elementy wykonane samodzielnie od gotowych komponentów.

\section{Instrukcja uruchomienia}
Podaj minimalny, jednoznaczny zestaw kroków potrzebnych do uruchomienia, zbudowania albo użycia rozwiązania.

% Kod przechowuj w katalogu code/ odpowiedniego rozdziału i wczytuj np.:
% \lstinputlisting[style=project-python,caption={Opis kodu},label={lst:przyklad}]{chapters/04-realizacja/code/program.py}

% TREŚĆ ROZDZIAŁU: zastąp wskazówki opisem testów i uzyskanych wyników.
% Każdą tabelę i każdy rysunek omów w tekście. Podawaj jednostki, warunki
% pomiaru oraz odwołanie do wymagania, które dany test weryfikuje.
\chapter{Weryfikacja i wyniki}
\label{chap:wyniki}

\section{Metodyka testów}
Opisz stanowisko, dane, warunki, procedurę i miary. Informacje powinny wystarczyć do powtórzenia najważniejszych testów.

\section{Wyniki}
Przedstaw wyniki w tabelach lub na wykresach, podając jednostki i warunki pomiaru. Nie powtarzaj w tekście wszystkich wartości widocznych w tabeli.

\section{Ocena spełnienia wymagań}
Dla każdego wymagania podaj wynik testu oraz status: spełnione, częściowo spełnione albo niespełnione. Odnieś się do kryteriów z \cref{chap:cel-wymagania}.

\section{Dyskusja}
Zinterpretuj rezultaty, wyjaśnij odstępstwa i wskaż ograniczenia wiarygodności oceny.

% TREŚĆ ROZDZIAŁU: zastąp wskazówki opisem przebiegu własnego projektu.
% Cały rozdział można wyłączyć w main.tex, jeśli prowadzący go nie wymaga.
\chapter{Organizacja i przebieg projektu}
\label{chap:organizacja}

\section{Etapy i harmonogram}
Porównaj planowane etapy z rzeczywistym przebiegiem. Wyjaśnij istotne zmiany zakresu lub terminów.

\section{Podział zadań}
W projekcie zespołowym jednoznacznie opisz odpowiedzialność poszczególnych osób. W projekcie indywidualnym wskaż własny wkład i wykorzystane elementy zewnętrzne.

\section{Zarządzanie wersjami i materiałami}
Krótko opisz repozytorium, zasady nazywania wersji, strukturę plików oraz sposób przechowywania danych i dokumentacji.

% TREŚĆ ROZDZIAŁU: zastąp wskazówki własnym podsumowaniem.
% Nie wprowadzaj tutaj nowych wyników, metod ani źródeł. Odwołuj się do dowodów
% przedstawionych wcześniej i jednoznacznie oceń realizację celu projektu.
\chapter{Podsumowanie i wnioski}
\label{chap:podsumowanie}

\section{Ocena realizacji celu}
Stwierdź, czy cel główny i cele szczegółowe osiągnięto, wskazując konkretne dowody z rozdziału wynikowego.

\section{Najważniejsze rezultaty}
Zbierz najważniejsze wyniki i wnioski bez wprowadzania nowych danych ani metod.

\section{Ograniczenia i dalszy rozwój}
Wskaż ograniczenia rozwiązania oraz konkretne, uzasadnione kierunki dalszych prac.


% === OŚWIADCZENIA KOŃCOWE — DOMYŚLNIE WŁĄCZONE ===
% Widoczność obu części ustawia się przez true/false w config/metadata.tex.
% Treść należy uzupełnić w odpowiednich plikach katalogu backmatter/.
\ifdefstring{\showaideclaration}{true}{% ============================================================================
% INFORMACJA O WYKORZYSTANIU NARZĘDZI AI
% ============================================================================
% Ta część realizuje wzór z załącznika nr 2A do Zasad korzystania z narzędzi
% AI przez studentów i doktorantów Politechniki Opolskiej.
%
% Wybierz dokładnie jeden wariant:
%   1) pozostaw treść poniżej, jeżeli AI wpłynęła na treść, strukturę,
%      rozwiązanie, kod, obliczenia, analizę, interpretację, wnioski lub inny
%      element materiału przekazywanego do oceny;
%   2) zastąp ją wariantem umieszczonym na końcu pliku, jeżeli nie korzystano
%      z AI w zakresie wymagającym ujawnienia.
%
% Zwykła autokorekta, sprawdzanie pisowni i inne funkcje techniczne, które nie
% wpływają na treść merytoryczną, co do zasady nie wymagają ujawnienia, chyba że
% prowadzący określi inaczej. Zakres użycia AI musi być zgodny z zasadami
% ustalonymi przez prowadzącego. Na jego żądanie trzeba przedstawić bardziej
% szczegółowe informacje, np. przykładowe polecenia, wyniki lub materiały
% robocze.
% ============================================================================

\chapter{Wykorzystanie narzędzi sztucznej inteligencji}
\label{chap:wykorzystanie-ai}

% === WARIANT 1: KORZYSTANO Z AI W ZAKRESIE WYMAGAJĄCYM UJAWNIENIA ===
% Dostosuj formy „korzystałem/am” i „przeze mnie” albo zastosuj formy mnogie,
% jeżeli projekt przygotował zespół.

W przygotowaniu niniejszego projektu korzystaliśmy z następujących narzędzi AI:

\noindent\textbf{Nazwa narzędzia lub narzędzi:}\par
Nazwa narzędzia, model lub wersja.

\medskip
\noindent\textbf{Cel i zakres wykorzystania:}\par
Opisz cel i zakres użycia, np. tworzenie pomysłów, analiza kodu, korekta
językowa albo proponowanie przypadków testowych.

\medskip
\noindent\textbf{Elementy przygotowane lub zmienione z wykorzystaniem AI:}\par
Wskaż konkretne rozdziały, fragmenty tekstu, kodu, obliczeń, analiz,
wizualizacji lub inne elementy projektu.

Wyniki uzyskane przy użyciu AI zostały przez nas sprawdzone, a za ostateczną
treść i poprawność przygotowanego materiału ponosimy odpowiedzialność.

% Nie wolno wprowadzać do niezatwierdzonych przez Uczelnię narzędzi AI m.in.
% danych osobowych, informacji poufnych i materiałów niepublicznych,
% nieopublikowanych danych i wyników badań, materiałów egzaminacyjnych, prac
% innych studentów, haseł, tokenów ani kluczy dostępu. Należy także przestrzegać
% praw autorskich, licencji, prawa do wizerunku i warunków użycia narzędzi.

% === WARIANT 2: NIE KORZYSTANO Z AI W ZAKRESIE WYMAGAJĄCYM UJAWNIENIA ===
% Aby użyć wariantu 2, usuń lub zakomentuj cały wariant 1 powyżej, a następnie
% odkomentuj poniższy akapit. Dostosuj formę osobową do autora albo zespołu.
%
% Oświadczamy, że podczas przygotowywania niniejszego projektu nie
% korzystaliśmy z narzędzi AI w sposób, który wpłynął na treść, strukturę,
% rozwiązanie, kod, obliczenia, analizę, interpretację, wnioski lub inne
% elementy materiału przekazywanego do oceny.
}{}
\ifdefstring{\showcontributions}{true}{% ============================================================================
% PODZIAŁ ZADAŃ I WKŁAD AUTORÓW
% ============================================================================
% Wymień konkretne rezultaty, nie tylko ogólne czynności. Udziały wszystkich
% autorów muszą sumować się do 100%. W projekcie indywidualnym wpisz jednego
% autora z udziałem 100%.
% ============================================================================

\chapter{Podział zadań i wkład autorów}
\label{chap:wklad-autorow}

Podział odpowiedzialności, wykonanych zadań i rezultatów przedstawiono
w \cref{tab:wklad-autorow}. Udział procentowy odnosi się do całego projektu,
a nie wyłącznie do objętości przygotowanego tekstu.

\begin{table}[htbp]
  \centering
  \caption{Podział zadań i udział autorów w realizacji projektu}
  \label{tab:wklad-autorow}
  \small
  \begin{tabularx}{\textwidth}{@{}p{2.8cm}X X p{1.5cm}@{}}
    \toprule
    Autor & Zadania i odpowiedzialność & Wytworzone rezultaty & Udział \\
    \midrule
    Imię i nazwisko
      & Np. opracowanie koncepcji, implementacja wybranego modułu i testy
      & Np. diagram architektury, kod modułu i przypadki testowe
      & \qty{100}{\percent} \\
    % Dla kolejnych autorów dodaj wiersze, a udziały zmień tak, aby suma
    % wynosiła 100%, np. 45% i 55%:
    % Imię i nazwisko & Zadania & Rezultaty & \qty{45}{\percent} \\
    \bottomrule
  \end{tabularx}
\end{table}

Autorzy potwierdzają zgodność powyższego zestawienia z rzeczywistym podziałem
pracy oraz akceptują przypisane zakresy odpowiedzialności i udziały procentowe.

% Historia repozytorium może potwierdzać przebieg pracy, ale sama liczba
% commitów nie jest miarą wartości wkładu poszczególnych osób.
}{}

% === BIBLIOGRAFIA — NIE USUWAJ, JEŻELI W RAPORCIE SĄ CYTOWANIA ===
% Źródła dodaje się do bibliografia.bib i przywołuje w tekście przez
% \cite{klucz}. Bibliografia zostanie utworzona automatycznie.
\printprojectbibliography

% === ZAŁĄCZNIKI — SEKCJA OPCJONALNA ===
% Polecenie \appendix przełącza numerację kolejnych rozdziałów na litery.
% Aby dodać załącznik, utwórz odpowiedni plik i usuń znak % przed \include.
% Jeżeli załączników nie ma, pozostaw przykładowe polecenie zakomentowane.
\appendix
% \include{chapters/08-zalacznik-a}

% === KONIEC DOKUMENTU — NICZEGO NIE DODAWAJ PONIŻEJ ===
\end{document}
